\documentclass[12pt,letterpaper]{article}

\usepackage[utf8]{inputenc}
\usepackage[T1]{fontenc}
\usepackage[spanish,es-tabla]{babel}
\usepackage{geometry}
\geometry{margin=2.5cm}

\usepackage{graphicx}
\usepackage{float}
\usepackage{booktabs}
\usepackage{array}
\usepackage{xcolor}
\usepackage{listings}
\usepackage{fancyhdr}
\usepackage{titlesec}
\usepackage{hyperref}

\definecolor{uabcGreen}{RGB}{0, 102, 51}
\definecolor{uabcGold}{RGB}{184, 134, 11}
\definecolor{codeBg}{RGB}{248, 249, 250}
\definecolor{codeBorder}{RGB}{210, 215, 220}
\definecolor{codeKeyword}{RGB}{0, 82, 164}
\definecolor{codeComment}{RGB}{106, 115, 125}
\definecolor{codeString}{RGB}{34, 134, 58}

\hypersetup{
    colorlinks=true,
    linkcolor=uabcGreen,
    urlcolor=codeKeyword,
    pdftitle={Arquitectura de Videojuego RPG por Turnos con Patrones GoF},
    pdfauthor={Adad Arias Palma}
}

\lstset{
    basicstyle=\ttfamily\footnotesize,
    backgroundcolor=\color{codeBg},
    frame=single,
    rulecolor=\color{codeBorder},
    breaklines=true,
    breakatwhitespace=false,
    numbers=left,
    numberstyle=\tiny\color{codeComment},
    stepnumber=1,
    numbersep=8pt,
    tabsize=4,
    showstringspaces=false,
    keywordstyle=\color{codeKeyword}\bfseries,
    commentstyle=\color{codeComment}\itshape,
    stringstyle=\color{codeString},
    literate=
        {á}{{\'a}}1 {é}{{\'e}}1 {í}{{\'i}}1 {ó}{{\'o}}1 {ú}{{\'u}}1
        {Á}{{\'A}}1 {É}{{\'E}}1 {Í}{{\'I}}1 {Ó}{{\'O}}1 {Ú}{{\'U}}1
        {ñ}{{\~n}}1 {Ñ}{{\~N}}1 {¡}{{!`}}1 {¿}{{?`}}1
}

\pagestyle{fancy}
\fancyhf{}
\fancyhead[L]{\small\color{codeComment} Análisis y Diseño de Sistemas -- Patrones GoF en Java}
\fancyhead[R]{\small\color{codeComment} UABC -- Facultad de Ingeniería}
\fancyfoot[C]{\thepage}
\renewcommand{\headrulewidth}{0.4pt}

\titleformat{\section}{\Large\bfseries\color{uabcGreen}}{\thesection.}{0.5em}{}
\titleformat{\subsection}{\large\bfseries\color{black}}{\thesubsection}{0.5em}{}

\begin{document}

% ==============================================================================
% PORTADA INSTITUCIONAL
% ==============================================================================
\begin{titlepage}
    \centering
    \vspace*{0.8cm}

    {\LARGE \bfseries \color{uabcGreen} UNIVERSIDAD AUTÓNOMA DE BAJA CALIFORNIA \par}
    \vspace{0.3cm}
    {\Large \bfseries FACULTAD DE INGENIERÍA \par}
    \vspace{0.2cm}
    {\large Programa de Ingeniería en Computación / Software \par}

    \vspace{1.2cm}
    \rule{\textwidth}{1.5pt}
    \vspace{0.6cm}

    {\Large \textbf{Reporte Técnico de Arquitectura de Software} \par}
    \vspace{0.4cm}
    {\huge \bfseries Diseño e Implementación de un Videojuego RPG por Turnos Aplicando Patrones de Diseño GoF \par}
    \vspace{0.4cm}
    {\normalsize \textit{Factory Method, Strategy, Decorator y Observer en Java 17 / Maven (Apache NetBeans 23)} \par}

    \vspace{0.6cm}
    \rule{\textwidth}{1.5pt}

    \vspace{1.8cm}
    \begin{tabular}{rl}
        \textbf{Asignatura:}  & Análisis y Diseño de Sistemas \\[0.3cm]
        \textbf{Alumno:}      & Adad Arias Palma \\[0.3cm]
        \textbf{Matrícula:}   & 1196869 \\[0.3cm]
        \textbf{Profesor:}    & Jose Martin Olguin Espinoza \\[0.3cm]
        \textbf{Ciclo:}       & 2026-2 \\[0.3cm]
        \textbf{Repositorio:} & \url{https://github.com/AdadAriasPalma/rpg-patrones-gof} \\
    \end{tabular}

    \vfill
    {\large Mexicali, Baja California, México \par}
    {\normalsize Ciclo Escolar 2026-2 \par}
\end{titlepage}

\tableofcontents
\newpage

% ==============================================================================
% SECCIÓN 1: CONTEXTO Y ARQUITECTURA GENERAL
% ==============================================================================
\section{Contexto y Arquitectura General}

El presente proyecto implementa desde cero un videojuego de rol (RPG) por turnos en consola bajo el estándar \textbf{Apache Maven} (\texttt{com.rpg.gof:rpg-patrones-gof:1.0.0}), compilado con Java 17, codificación UTF-8 y configurado con \texttt{exec-maven-plugin} (\texttt{3.1.0}) para su ejecución inmediata en \textbf{Apache NetBeans IDE 23} o mediante línea de comandos. El código fuente completo y sus artefactos se encuentran alojados en el repositorio público de GitHub: \url{https://github.com/AdadAriasPalma/rpg-patrones-gof}.

El objetivo arquitectónico central consiste en evitar el acoplamiento rígido típico de las jerarquías monolíticas en videojuegos (explosión de subclases por combinaciones de clase de héroe, tipo de ataque y equipamiento) y separar por completo la lógica determinista de combate respecto al mecanismo de registro o interfaz de salida. Para ello, el sistema se organiza en siete paquetes cohesivos con responsabilidades claramente delimitadas:

\begin{table}[H]
    \centering
    \caption{Distribución modular por paquetes del proyecto \texttt{rpg-patrones-gof}.}
    \label{tab:paquetes}
    \begin{tabular}{>{\ttfamily}l l p{6.8cm}}
        \toprule
        \normalfont\textbf{Paquete} & \textbf{Rol Arquitectónico} & \textbf{Componentes Principales} \\
        \midrule
        com.rpg.gof.model     & Modelo de Dominio   & \texttt{Personaje}, \texttt{PersonajeBase}, \texttt{ResultadoAtaque}. \\
        com.rpg.gof.strategy  & Patrón Strategy     & \texttt{EstrategiaAtaque} y 6 estrategias concretas deterministas (\texttt{AtaqueEspada}, \texttt{AtaqueHechizo}, \texttt{AtaqueFlecha}, \texttt{AtaquePunalada}, \texttt{AtaqueGarrote}, \texttt{AtaqueAlientoFuego}). \\
        com.rpg.gof.factory   & Patrón Factory Method & \texttt{CreadorPersonaje} (Guerrero, Mago, Arquero, Pícaro) y \texttt{CreadorEnemigo} (Goblin, Orco, Dragón). \\
        com.rpg.gof.decorator & Patrón Decorator    & \texttt{PersonajeDecorator} y 5 decoradores apilables (\texttt{Espada}, \texttt{Armadura}, \texttt{Veneno}, \texttt{AmuletoRegeneracion}, \texttt{Bendicion}). \\
        com.rpg.gof.observer  & Patrón Observer     & \texttt{ObservadorCombate} y \texttt{BitacoraCombate}. \\
        com.rpg.gof.engine    & Motor / Subject     & \texttt{MotorCombate}. \\
        com.rpg.gof           & Cliente / Orquestador & \texttt{Main} (ejecución de las 3 corridas de prueba). \\
        \bottomrule
    \end{tabular}
\end{table}

Todas las fórmulas de combate fueron diseñadas de forma determinista con el fin de garantizar la reproducibilidad exacta de los escenarios de prueba y facilitar la verificación matemática de cada turno.

% ==============================================================================
% SECCIÓN 2: JUSTIFICACIÓN DE PATRONES GOF Y PRINCIPIOS SOLID
% ==============================================================================
\section{Justificación de Patrones GoF y Principios SOLID}

\subsection{Patrón Strategy (\texttt{com.rpg.gof.strategy})}
\textbf{Problema que resuelve:} En un RPG, los distintos combatientes emplean mecánicas ofensivas heterogéneas (ataques cuerpo a cuerpo mitigados por armadura, hechizos arcanos que ignoran dos tercios de la defensa, flechas perforantes que ignoran la mitad, o puñaladas sigilosas que escalan con la velocidad). Codificar estas variantes mediante condicionales (\texttt{if/switch}) dentro de \texttt{PersonajeBase} violaría el principio Abierto/Cerrado.

\textbf{Solución aplicada:} Se define la interfaz \texttt{EstrategiaAtaque} con el contrato \texttt{atacar(Personaje atacante, Personaje objetivo)} retornando un objeto de valor inmutable \texttt{ResultadoAtaque}. Cada algoritmo reside en su propia clase concreta, permitiendo intercambiar la estrategia de un combatiente en tiempo de ejecución mediante \texttt{setEstrategia(...)}.

\subsection{Patrón Factory Method (\texttt{com.rpg.gof.factory})}
\textbf{Problema que resuelve:} Instanciar directamente \texttt{PersonajeBase} desde la capa cliente obligaría a conocer los valores exactos de HP, ataque, defensa, velocidad y la estrategia asociada a cada arquetipo de héroe o monstruo, dispersando reglas de balanceo por todo el código.

\textbf{Solución aplicada:} Se establecen dos jerarquías de creadores abstractos: \texttt{CreadorPersonaje} (con el método fábrica \texttt{crearPersonaje(String)} y el \textit{hook} \texttt{validarNombre(String)}) y \texttt{CreadorEnemigo} (con \texttt{crearEnemigo()}). Cada subclase concreta encapsula la configuración oficial de estadísticas y la estrategia inicial correspondiente.

\subsection{Patrón Decorator (\texttt{com.rpg.gof.decorator})}
\textbf{Problema que resuelve:} Si un héroe puede equipar combinaciones arbitrarias de espada, armadura, veneno, amuleto de regeneración y bendición, modelar esto con herencia generaría una explosión combinatoria ($2^5 = 32$ subclases por cada clase de héroe).

\textbf{Solución aplicada:} \texttt{PersonajeDecorator} implementa \texttt{Personaje} y envuelve a una instancia \texttt{protected final Personaje personaje}. Los decoradores concretos extienden esta base y modifican los atributos o comportamientos de forma acumulativa y en cualquier orden de anidamiento:
\begin{itemize}
    \item \texttt{EspadaDecorator}: incrementa el ataque en $+10$.
    \item \texttt{ArmaduraDecorator}: incrementa la defensa en $+8$.
    \item \texttt{VenenoDecorator}: añade $+5$ de ataque, suma $+5$ de daño al resultado de \texttt{atacar()} y concatena \texttt{" + veneno"} al efecto.
    \item \texttt{AmuletoRegeneracionDecorator}: aporta $+3$ de defensa y regenera $+2$ HP por turno sin superar \texttt{hpMax}.
    \item \texttt{BendicionDecorator}: multiplica ataque ($\times 1.15$ truncado), defensa ($\times 1.15$ redondeado) y velocidad ($\times 1.15$ truncado).
\end{itemize}

\textbf{Consideración crítica de diseño:} Al combinar \texttt{Decorator} con \texttt{Strategy}, si \texttt{PersonajeBase.atacar()} pasara \texttt{this} a la estrategia, la fórmula perdería los bonos de los decoradores externos. Para resolverlo con total independencia del orden de envoltura, \texttt{PersonajeDecorator} emplea un contexto \texttt{ThreadLocal<Personaje>} que conserva la referencia al decorador más externo mientras la llamada desciende hacia el componente base y atraviesa decoradores intermedios como \texttt{VenenoDecorator}.

\subsection{Patrón Observer (\texttt{com.rpg.gof.observer} y \texttt{com.rpg.gof.engine})}
\textbf{Problema que resuelve:} Acoplar el ciclo de combate a \texttt{System.out.println} impediría reutilizar el motor en una interfaz gráfica (Swing/JavaFX), un servidor web o pruebas automatizadas en memoria.

\textbf{Solución aplicada:} \texttt{MotorCombate} actúa como \textit{Subject} manteniendo una lista de \texttt{ObservadorCombate}. Durante el ciclo de turnos emite eventos estructurados (\texttt{onInicioCombate}, \texttt{onTurnoEjecutado}, \texttt{onGolpeCritico}, \texttt{onRegeneracion} y \texttt{onFinCombate}). \texttt{BitacoraCombate} implementa dicha interfaz para formatear la salida en consola y simultáneamente conservar el historial en memoria.

\subsection{Alineación con Principios SOLID}

\begin{table}[H]
    \centering
    \caption{Correspondencia entre los principios SOLID y las decisiones arquitectónicas del proyecto.}
    \begin{tabular}{l p{10.5cm}}
        \toprule
        \textbf{Principio} & \textbf{Evidencia en la Arquitectura} \\
        \midrule
        \textbf{S} -- Responsabilidad Única & Cada clase cumple un único propósito: \texttt{MotorCombate} solo orquesta turnos, \texttt{BitacoraCombate} solo registra eventos, cada \texttt{EstrategiaAtaque} encapsula una sola fórmula y cada \texttt{Creador} instancia un único arquetipo. \\
        \textbf{O} -- Abierto / Cerrado & Pueden agregarse nuevas clases de héroe, enemigos, estrategias, decoradores u observadores sin modificar las clases existentes ni el núcleo de \texttt{MotorCombate}. \\
        \textbf{L} -- Sustitución de Liskov & Todo decorador concreto y \texttt{PersonajeBase} implementan íntegramente el contrato \texttt{Personaje}, siendo completamente intercambiables ante \texttt{MotorCombate}. \\
        \textbf{I} -- Segregación de Interfaces & Se definen interfaces focalizadas y cohesivas (\texttt{EstrategiaAtaque} con un método principal, \texttt{ObservadorCombate} orientada exclusivamente a eventos de combate). \\
        \textbf{D} -- Inversión de Dependencias & \texttt{MotorCombate} depende de las abstracciones \texttt{Personaje} y \texttt{ObservadorCombate}; \texttt{PersonajeBase} depende de la abstracción \texttt{EstrategiaAtaque}; \texttt{Main} depende de las fábricas abstractas \texttt{CreadorPersonaje} y \texttt{CreadorEnemigo}. \\
        \bottomrule
    \end{tabular}
\end{table}

% ==============================================================================
% SECCIÓN 3: DIAGRAMAS UML
% ==============================================================================
\section{Diagramas UML}

Los diagramas fueron modelados formalmente en \textbf{PlantUML} dentro del directorio \texttt{diagramas/} para documentar tanto la vista estructural (clases, interfaces y dependencias entre los cuatro patrones) como la vista dinámica (secuencia completa de creación, decoración y ejecución de combate).

\subsection{Diagrama de Clases}

\begin{figure}[H]
    \centering
    \IfFileExists{diagramas/diagrama_clases.png}{
        \includegraphics[width=\textwidth]{diagramas/diagrama_clases.png}
    }{
        \fbox{\parbox{0.92\textwidth}{\centering\vspace{1cm}
            \textbf{[Figura reservada: \texttt{diagramas/diagrama\_clases.png}]}\\[0.25cm]
            \small Para incluir la imagen compilada, renderice el archivo \texttt{diagramas/diagrama\_clases.puml} con PlantUML y guárdelo en la ruta indicada.
            \vspace{1cm}}}
    }
    \caption{Diagrama de clases completo integrando los patrones Factory Method, Strategy, Decorator y Observer.}
    \label{fig:diagrama_clases}
\end{figure}

\begin{lstlisting}[caption={Código fuente PlantUML del Diagrama de Clases (\texttt{diagramas/diagrama\_clases.puml}).}]
@startuml
title Diagrama de Clases - Videojuego RPG por Turnos (Patrones GoF)

skinparam packageStyle rectangle
skinparam classAttributeIconSize 0
skinparam linetype ortho
skinparam shadowing false

package "com.rpg.gof.model" {
    interface Personaje <<Component>> {
        + getNombre() : String
        + getAtaque() : int
        + getDefensa() : int
        + getVelocidad() : int
        + getHp() : int
        + getHpMax() : int
        + setHp(hp : int) : void
        + getEstrategia() : EstrategiaAtaque
        + setEstrategia(e : EstrategiaAtaque) : void
        + atacar(objetivo : Personaje) : ResultadoAtaque
        + getDescripcion() : String
        + regenerarPorTurno() : int
        + estaVivo() : boolean
    }

    class PersonajeBase <<ConcreteComponent>> {
        - nombre : String
        - hp : int
        - hpMax : int
        - ataque : int
        - defensa : int
        - velocidad : int
        - estrategia : EstrategiaAtaque
        + PersonajeBase(nombre, hpMax, ataque, defensa, velocidad, estrategia)
        + atacar(objetivo : Personaje) : ResultadoAtaque
        + getDescripcion() : String
    }

    class ResultadoAtaque <<ValueObject>> {
        - dano : int
        - critico : boolean
        - efecto : String
        + ResultadoAtaque(dano : int, critico : boolean, efecto : String)
        + getDano() : int
        + isCritico() : boolean
        + getEfecto() : String
    }
}

package "com.rpg.gof.strategy" {
    interface EstrategiaAtaque <<Strategy>> {
        + atacar(atacante : Personaje, objetivo : Personaje) : ResultadoAtaque
        + getNombreEstrategia() : String
    }
    class AtaqueEspada <<ConcreteStrategy>>
    class AtaqueHechizo <<ConcreteStrategy>>
    class AtaqueFlecha <<ConcreteStrategy>>
    class AtaquePunalada <<ConcreteStrategy>>
    class AtaqueGarrote <<ConcreteStrategy>>
    class AtaqueAlientoFuego <<ConcreteStrategy>>
}

package "com.rpg.gof.decorator" {
    abstract class PersonajeDecorator <<Decorator>> {
        # personaje : Personaje
        # PersonajeDecorator(personaje : Personaje)
        + atacar(objetivo : Personaje) : ResultadoAtaque
        + regenerarPorTurno() : int
    }
    class EspadaDecorator <<ConcreteDecorator>>
    class ArmaduraDecorator <<ConcreteDecorator>>
    class VenenoDecorator <<ConcreteDecorator>>
    class AmuletoRegeneracionDecorator <<ConcreteDecorator>>
    class BendicionDecorator <<ConcreteDecorator>>
}

package "com.rpg.gof.factory" {
    abstract class CreadorPersonaje <<Creator>> {
        + {abstract} crearPersonaje(nombre : String) : Personaje
        # validarNombre(nombre : String) : void
    }
    class CreadorGuerrero <<ConcreteCreator>>
    class CreadorMago <<ConcreteCreator>>
    class CreadorArquero <<ConcreteCreator>>
    class CreadorPicaro <<ConcreteCreator>>

    abstract class CreadorEnemigo <<Creator>> {
        + {abstract} crearEnemigo() : Personaje
    }
    class CreadorGoblin <<ConcreteCreator>>
    class CreadorOrco <<ConcreteCreator>>
    class CreadorDragon <<ConcreteCreator>>
}

package "com.rpg.gof.observer" {
    interface ObservadorCombate <<Observer>> {
        + onInicioCombate(heroe : Personaje, enemigo : Personaje) : void
        + onTurnoEjecutado(numeroTurno : int, atacante : Personaje, defensor : Personaje, resultado : ResultadoAtaque) : void
        + onGolpeCritico(atacante : Personaje, defensor : Personaje, resultado : ResultadoAtaque) : void
        + onRegeneracion(personaje : Personaje, puntosRegenerados : int, hpActual : int, hpMax : int) : void
        + onFinCombate(ganador : Personaje, perdedor : Personaje, totalTurnos : int) : void
    }
    class BitacoraCombate <<ConcreteObserver>> {
        - historial : List<String>
        + getHistorial() : List<String>
    }
}

package "com.rpg.gof.engine" {
    class MotorCombate <<Subject>> {
        - observadores : List<ObservadorCombate>
        + agregarObservador(observador : ObservadorCombate) : void
        + eliminarObservador(observador : ObservadorCombate) : void
        + iniciarCombate(heroe : Personaje, enemigo : Personaje) : Personaje
    }
}

PersonajeBase ..|> Personaje
PersonajeBase o--> "1" EstrategiaAtaque : estrategia
PersonajeDecorator ..|> Personaje
PersonajeDecorator o--> "1" Personaje : envuelve
EspadaDecorator --|> PersonajeDecorator
ArmaduraDecorator --|> PersonajeDecorator
VenenoDecorator --|> PersonajeDecorator
AmuletoRegeneracionDecorator --|> PersonajeDecorator
BendicionDecorator --|> PersonajeDecorator
BitacoraCombate ..|> ObservadorCombate
MotorCombate o--> "*" ObservadorCombate : notifica
@enduml
\end{lstlisting}

\subsection{Diagrama de Secuencia}

\begin{figure}[H]
    \centering
    \IfFileExists{diagramas/diagrama_secuencia.png}{
        \includegraphics[width=\textwidth]{diagramas/diagrama_secuencia.png}
    }{
        \fbox{\parbox{0.92\textwidth}{\centering\vspace{1cm}
            \textbf{[Figura reservada: \texttt{diagramas/diagrama\_secuencia.png}]}\\[0.25cm]
            \small Para incluir la imagen compilada, renderice el archivo \texttt{diagramas/diagrama\_secuencia.puml} con PlantUML y guárdelo en la ruta indicada.
            \vspace{1cm}}}
    }
    \caption{Diagrama de secuencia del flujo de creación, decoración en cadena y ejecución de turno.}
    \label{fig:diagrama_secuencia}
\end{figure}

\begin{lstlisting}[caption={Código fuente PlantUML del Diagrama de Secuencia (\texttt{diagramas/diagrama\_secuencia.puml}).}]
@startuml
title Diagrama de Secuencia - Creación (Factory), Equipamiento (Decorator) y Turno de Combate (Strategy + Observer)

skinparam shadowing false
skinparam sequenceMessageAlign center

actor Usuario
participant "Main" as Main
participant "fabHeroe : CreadorPicaro" as FabHeroe
participant "fabEnemigo : CreadorDragon" as FabEnemigo
participant "sombraBase : PersonajeBase" as Base
participant "espada : EspadaDecorator" as Espada
participant "veneno : VenenoDecorator" as Veneno
participant "sombra : BendicionDecorator" as Bendicion
participant "estrategia : AtaquePunalada" as Estrategia
participant "motor : MotorCombate" as Motor
participant "bitacora : BitacoraCombate" as Bitacora

== Fase 1: Creación mediante Factory Method ==
Usuario -> Main : main(args)
Main -> FabHeroe : crearPersonaje("Sombra")
activate FabHeroe
FabHeroe -> FabHeroe : validarNombre("Sombra")
FabHeroe -> Base ** : new PersonajeBase("Sombra", 105, 26, 12, 22, AtaquePunalada)
FabHeroe --> Main : sombraBase
deactivate FabHeroe

Main -> FabEnemigo : crearEnemigo()
activate FabEnemigo
FabEnemigo --> Main : dragon ("Dragón ancestral")
deactivate FabEnemigo

== Fase 2: Equipamiento Dinámico en Cadena (Decorator) ==
Main -> Espada ** : new EspadaDecorator(sombraBase)
Main -> Veneno ** : new VenenoDecorator(espada)
Main -> Bendicion ** : new BendicionDecorator(veneno)

== Fase 3: Registro de Observador e Inicio de Combate (Observer) ==
Main -> Motor : agregarObservador(bitacora)
Main -> Motor : iniciarCombate(sombra, dragon)
activate Motor
Motor -> Bitacora : onInicioCombate(sombra, dragon)
activate Bitacora
Bitacora --> Motor : imprime fichas iniciales e iniciativa
deactivate Bitacora

== Fase 4: Ejecución de Turno (Decorator + Strategy + Observer) ==
loop mientras ambos personajes estén vivos
    Motor -> Bendicion : atacar(dragon)
    activate Bendicion
    Bendicion -> Veneno : atacar(dragon)
    activate Veneno
    Veneno -> Espada : atacar(dragon)
    activate Espada
    Espada -> Base : getEstrategia()
    Base --> Espada : estrategia
    Espada -> Estrategia : atacar(sombra, dragon)
    activate Estrategia
    Estrategia -> Bendicion : getAtaque() / getVelocidad()
    Bendicion --> Estrategia : ATK=47, VEL=25 (>=20 -> Crítico)
    Estrategia --> Espada : ResultadoAtaque(43, true, "puñalada furtiva")
    deactivate Estrategia
    Espada --> Veneno : ResultadoAtaque(43, true, "puñalada furtiva")
    deactivate Espada
    Veneno --> Bendicion : ResultadoAtaque(48, true, "puñalada furtiva + veneno")
    deactivate Veneno
    Bendicion --> Motor : ResultadoAtaque(48, true, "puñalada furtiva + veneno")
    deactivate Bendicion

    Motor -> Motor : dragon.setHp(hp - 48)
    Motor -> Bitacora : onGolpeCritico(sombra, dragon, resultado)
    Motor -> Bitacora : onTurnoEjecutado(turno, sombra, dragon, resultado)
    Motor -> Bendicion : regenerarPorTurno()
    Bendicion --> Motor : 0 (sin amuleto)
end

Motor -> Bitacora : onFinCombate(sombra, dragon, totalTurnos)
activate Bitacora
Bitacora --> Motor : imprime resumen y banner ¡VICTORIA PARA SOMBRA!
deactivate Bitacora

Motor --> Main : ganador (sombra)
deactivate Motor
@enduml
\end{lstlisting}

% ==============================================================================
% SECCIÓN 4: FRAGMENTOS CLAVE DE IMPLEMENTACIÓN
% ==============================================================================
\section{Fragmentos Clave de Implementación}

A continuación se presentan los fragmentos representativos que materializan los contratos e interacciones críticas de cada patrón en Java 17.

\subsection{Strategy: \texttt{AtaquePunalada.java}}
\begin{lstlisting}[language=Java, caption={Estrategia determinista con evaluación de golpe crítico por velocidad.}]
public class AtaquePunalada implements EstrategiaAtaque {

    private static final int UMBRAL_VELOCIDAD_CRITICA = 20;

    @Override
    public ResultadoAtaque atacar(Personaje atacante, Personaje objetivo) {
        int danoBase = Math.max(1, atacante.getAtaque() - objetivo.getDefensa());
        boolean critico = atacante.getVelocidad() >= UMBRAL_VELOCIDAD_CRITICA;
        int danoFinal = critico ? Math.max(1, (int) (danoBase * 1.5)) : danoBase;
        return new ResultadoAtaque(danoFinal, critico, "puñalada furtiva");
    }
}
\end{lstlisting}

\subsection{Factory Method: \texttt{CreadorPersonaje.java} y \texttt{CreadorPicaro.java}}
\begin{lstlisting}[language=Java, caption={Creador abstracto con hook de validación y fábrica concreta de Pícaro.}]
public abstract class CreadorPersonaje {
    public abstract Personaje crearPersonaje(String nombre);

    protected void validarNombre(String nombre) {
        if (nombre == null || nombre.trim().isEmpty()) {
            throw new IllegalArgumentException("El nombre del héroe no puede ser nulo o vacío.");
        }
    }
}

public class CreadorPicaro extends CreadorPersonaje {
    @Override
    public Personaje crearPersonaje(String nombre) {
        validarNombre(nombre);
        return new PersonajeBase(nombre.trim(), 105, 26, 12, 22, new AtaquePunalada());
    }
}
\end{lstlisting}

\subsection{Decorator: \texttt{PersonajeDecorator.java} y \texttt{VenenoDecorator.java}}
\begin{lstlisting}[language=Java, caption={Propagación del decorador externo hacia la estrategia y adición de daño por veneno.}]
public abstract class PersonajeDecorator implements Personaje {
    private static final ThreadLocal<Personaje> ATACANTE_DECORADO = new ThreadLocal<>();
    protected final Personaje personaje;

    @Override
    public ResultadoAtaque atacar(Personaje objetivo) {
        Objects.requireNonNull(objetivo, "El objetivo del ataque no puede ser nulo.");
        boolean esExterior = (ATACANTE_DECORADO.get() == null);
        if (esExterior) {
            ATACANTE_DECORADO.set(this);
        }
        try {
            if (personaje instanceof PersonajeDecorator) {
                return personaje.atacar(objetivo);
            }
            return personaje.getEstrategia().atacar(ATACANTE_DECORADO.get(), objetivo);
        } finally {
            if (esExterior) {
                ATACANTE_DECORADO.remove();
            }
        }
    }
}

public class VenenoDecorator extends PersonajeDecorator {
    @Override
    public int getAtaque() {
        return personaje.getAtaque() + 5;
    }

    @Override
    public ResultadoAtaque atacar(Personaje objetivo) {
        ResultadoAtaque resultadoBase = super.atacar(objetivo);
        return new ResultadoAtaque(
                resultadoBase.getDano() + 5,
                resultadoBase.isCritico(),
                resultadoBase.getEfecto() + " + veneno"
        );
    }
}
\end{lstlisting}

\subsection{Observer y Motor: \texttt{MotorCombate.java}}
\begin{lstlisting}[language=Java, caption={Orquestación de turnos y notificación desacoplada a observadores.}]
public Personaje iniciarCombate(Personaje heroe, Personaje enemigo) {
    notificarInicioCombate(heroe, enemigo);

    Personaje primero = heroe.getVelocidad() >= enemigo.getVelocidad() ? heroe : enemigo;
    Personaje segundo = (primero == heroe) ? enemigo : heroe;

    int numeroTurno = 0;
    while (heroe.estaVivo() && enemigo.estaVivo()) {
        numeroTurno++;
        ejecutarAccion(numeroTurno, primero, segundo);
        if (segundo.estaVivo()) {
            ejecutarAccion(numeroTurno, segundo, primero);
        }
    }

    Personaje ganador = heroe.estaVivo() ? heroe : enemigo;
    Personaje perdedor = (ganador == heroe) ? enemigo : heroe;
    notificarFinCombate(ganador, perdedor, numeroTurno);
    return ganador;
}
\end{lstlisting}

% ==============================================================================
% SECCIÓN 5: EVIDENCIA DE EJECUCIÓN
% ==============================================================================
\section{Evidencia de Ejecución}

El proyecto se compiló y ejecutó mediante Maven (\texttt{mvn clean compile exec:java}) generando el archivo \texttt{EJECUCION.txt}. La Tabla~\ref{tab:resumen_corridas} resume los resultados deterministas verificados en las tres corridas.

\begin{table}[H]
    \centering
    \caption{Resumen cuantitativo de las tres corridas ejecutadas en \texttt{Main.java}.}
    \label{tab:resumen_corridas}
    \begin{tabular}{llccccc}
        \toprule
        \textbf{Corrida} & \textbf{Enfrentamiento} & \textbf{Turnos} & \textbf{Daño Total} & \textbf{Críticos} & \textbf{HP Regen.} & \textbf{Ganador (HP)} \\
        \midrule
        1 & Aria vs Goblin salvaje     & 2 & 91  & 0 & 0  & Aria (99/100) \\
        2 & Kael vs Orco brutal        & 6 & 168 & 0 & 10 & Kael (120/140) \\
        3 & Sombra vs Dragón ancestral & 4 & 273 & 4 & 0  & Sombra (24/105) \\
        \bottomrule
    \end{tabular}
\end{table}

\subsection{Salida Real Capturada (\texttt{EJECUCION.txt})}

\begin{lstlisting}[caption={Salida íntegra generada al compilar y ejecutar el proyecto con Maven.}]
[INFO] --- compiler:3.13.0:compile (default-compile) @ rpg-patrones-gof ---
[INFO] Recompiling the module because of changed source code.
[INFO] Compiling 29 source files with javac [debug target 17] to target\classes
[INFO] --- exec:3.1.0:java (default-cli) @ rpg-patrones-gof ---

################################################################################
CORRIDA 1: Aria (Mago) + Espada + Veneno + Bendición vs Goblin salvaje
################################################################################
================================================================================
                         INICIO DEL COMBATE POR TURNOS                          
================================================================================
HÉROE   | Aria (base) + Espada + Veneno + Bendición        | HP:100/100 | ATK:40 | DEF:12 | VEL:17 | Estrategia: AtaqueHechizo
ENEMIGO | Goblin salvaje (base)                            | HP: 60/ 60 | ATK:13 | DEF: 2 | VEL:14 | Estrategia: AtaquePunalada
--------------------------------------------------------------------------------
Orden de iniciativa por velocidad: Aria ataca primero.
--------------------------------------------------------------------------------
[Turno 01] Aria             ejecuta "hechizo arcano + veneno"    -> 45 de daño | Goblin salvaje HP: 15/60
[Turno 01] Goblin salvaje   ejecuta "puñalada furtiva"           ->  1 de daño | Aria HP: 99/100
[Turno 02] Aria             ejecuta "hechizo arcano + veneno"    -> 45 de daño | Goblin salvaje HP: 0/60
--------------------------------------------------------------------------------
RESUMEN DEL COMBATE:
  - Turnos transcurridos : 2
  - Daño total acumulado : 91
  - Golpes críticos      : 0
  - HP total regenerado  : 0
  - Estado del ganador   : Aria (base) + Espada + Veneno + Bendición (HP 99/100)
  - Estado del derrotado : Goblin salvaje (HP 0/60)
********************************************************************************
                 ¡VICTORIA PARA ARIA!
********************************************************************************

################################################################################
CORRIDA 2: Kael (Guerrero) + Espada + Armadura + Amuleto de Regeneración vs Orco brutal
################################################################################
================================================================================
                         INICIO DEL COMBATE POR TURNOS                          
================================================================================
HÉROE   | Kael (base) + Espada + Armadura + Amuleto de Regeneración | HP:140/140 | ATK:35 | DEF:26 | VEL:10 | Estrategia: AtaqueEspada
ENEMIGO | Orco brutal (base)                               | HP:130/130 | ATK:28 | DEF:12 | VEL: 8 | Estrategia: AtaqueGarrote
--------------------------------------------------------------------------------
Orden de iniciativa por velocidad: Kael ataca primero.
--------------------------------------------------------------------------------
[Turno 01] Kael             ejecuta "corte de espada"            -> 23 de daño | Orco brutal HP: 107/130
[Turno 01] Orco brutal      ejecuta "golpe de garrote"           ->  6 de daño | Kael HP: 134/140
[Turno 02] Kael             ejecuta "corte de espada"            -> 23 de daño | Orco brutal HP: 84/130
   >>> [REGENERACIÓN] Kael recupera +2 HP mediante su amuleto -> HP actual: 136/140
[Turno 02] Orco brutal      ejecuta "golpe de garrote"           ->  6 de daño | Kael HP: 130/140
[Turno 03] Kael             ejecuta "corte de espada"            -> 23 de daño | Orco brutal HP: 61/130
   >>> [REGENERACIÓN] Kael recupera +2 HP mediante su amuleto -> HP actual: 132/140
[Turno 03] Orco brutal      ejecuta "golpe de garrote"           ->  6 de daño | Kael HP: 126/140
[Turno 04] Kael             ejecuta "corte de espada"            -> 23 de daño | Orco brutal HP: 38/130
   >>> [REGENERACIÓN] Kael recupera +2 HP mediante su amuleto -> HP actual: 128/140
[Turno 04] Orco brutal      ejecuta "golpe de garrote"           ->  6 de daño | Kael HP: 122/140
[Turno 05] Kael             ejecuta "corte de espada"            -> 23 de daño | Orco brutal HP: 15/130
   >>> [REGENERACIÓN] Kael recupera +2 HP mediante su amuleto -> HP actual: 124/140
[Turno 05] Orco brutal      ejecuta "golpe de garrote"           ->  6 de daño | Kael HP: 118/140
[Turno 06] Kael             ejecuta "corte de espada"            -> 23 de daño | Orco brutal HP: 0/130
   >>> [REGENERACIÓN] Kael recupera +2 HP mediante su amuleto -> HP actual: 120/140
--------------------------------------------------------------------------------
RESUMEN DEL COMBATE:
  - Turnos transcurridos : 6
  - Daño total acumulado : 168
  - Golpes críticos      : 0
  - HP total regenerado  : 10
  - Estado del ganador   : Kael (base) + Espada + Armadura + Amuleto de Regeneración (HP 120/140)
  - Estado del derrotado : Orco brutal (HP 0/130)
********************************************************************************
                 ¡VICTORIA PARA KAEL!
********************************************************************************

################################################################################
CORRIDA 3: Sombra (Pícaro) + Espada + Veneno + Bendición vs Dragón ancestral
################################################################################
================================================================================
                         INICIO DEL COMBATE POR TURNOS                          
================================================================================
HÉROE   | Sombra (base) + Espada + Veneno + Bendición      | HP:105/105 | ATK:47 | DEF:14 | VEL:25 | Estrategia: AtaquePunalada
ENEMIGO | Dragón ancestral (base)                          | HP:180/180 | ATK:34 | DEF:18 | VEL:16 | Estrategia: AtaqueAlientoFuego
--------------------------------------------------------------------------------
Orden de iniciativa por velocidad: Sombra ataca primero.
--------------------------------------------------------------------------------
   >>> [¡CRÍTICO!] Sombra aprovecha su velocidad (25) para asestar un golpe letal sobre Dragón ancestral.
[Turno 01] Sombra           ejecuta "puñalada furtiva + veneno"  -> 48 de daño [CRÍTICO] | Dragón ancestral HP: 132/180
[Turno 01] Dragón ancestral ejecuta "aliento de fuego"           -> 27 de daño | Sombra HP: 78/105
   >>> [¡CRÍTICO!] Sombra aprovecha su velocidad (25) para asestar un golpe letal sobre Dragón ancestral.
[Turno 02] Sombra           ejecuta "puñalada furtiva + veneno"  -> 48 de daño [CRÍTICO] | Dragón ancestral HP: 84/180
[Turno 02] Dragón ancestral ejecuta "aliento de fuego"           -> 27 de daño | Sombra HP: 51/105
   >>> [¡CRÍTICO!] Sombra aprovecha su velocidad (25) para asestar un golpe letal sobre Dragón ancestral.
[Turno 03] Sombra           ejecuta "puñalada furtiva + veneno"  -> 48 de daño [CRÍTICO] | Dragón ancestral HP: 36/180
[Turno 03] Dragón ancestral ejecuta "aliento de fuego"           -> 27 de daño | Sombra HP: 24/105
   >>> [¡CRÍTICO!] Sombra aprovecha su velocidad (25) para asestar un golpe letal sobre Dragón ancestral.
[Turno 04] Sombra           ejecuta "puñalada furtiva + veneno"  -> 48 de daño [CRÍTICO] | Dragón ancestral HP: 0/180
--------------------------------------------------------------------------------
RESUMEN DEL COMBATE:
  - Turnos transcurridos : 4
  - Daño total acumulado : 273
  - Golpes críticos      : 4
  - HP total regenerado  : 0
  - Estado del ganador   : Sombra (base) + Espada + Veneno + Bendición (HP 24/105)
  - Estado del derrotado : Dragón ancestral (HP 0/180)
********************************************************************************
                 ¡VICTORIA PARA SOMBRA!
********************************************************************************
[INFO] BUILD SUCCESS
\end{lstlisting}

\end{document}
